\section{Conclusion}

Base-\(B\) multiplication can be viewed as bilinear projection followed by
nonlinear carry flow. The outer-product matrix records pairwise digit
interactions, anti-diagonal projection aggregates those interactions into raw
degree coefficients, and carry normalization transports excess mass along the
degree axis while preserving numerical value.

Carry is locally activated, globally propagated, and geometrically
conditioned. It is locally activated by the threshold \(c_k+r_k\ge B\), globally
propagated because outgoing carry changes the next degree, and geometrically
conditioned because support-sumset representation counts control where raw mass
can concentrate. This gives a coherent language for studying multiplication as
projection followed by normalization, and for relating carry complexity to the
geometry of digit supports.
